% Section 7.3, from lectures/L07/README.md: what you should be able to do afterwards, the
% questions to test yourself, and the next lecture.

\section{Review}
\label{c7:sec:review}

\needspace{6\baselineskip}
\subsection*{What you should be able to do now}
After this chapter, you should be able to:
\begin{itemize}
  \item Name five sources of concurrency in a kernel and say which of them a single-CPU kernel has.
  \item Say why making every field of a structure atomic does not make the structure thread-safe.
  \item Choose between a spinlock and a mutex, and give the one-sentence reason.
  \item Say what \code{spin_lock_irqsave} saves and restores, and what goes wrong without it.
  \item Recognise atomic context in unfamiliar code, and list what may not be done there.
  \item Read a lockdep report and say which two locks, taken in which orders, it is complaining
    about.
  \item Explain why lockdep can report a deadlock that has not happened.
\end{itemize}

\needspace{6\baselineskip}
\subsection*{Questions to test yourself}
\begin{enumerate}
  \item Your driver works perfectly on a single-CPU target and corrupts data on a dual-CPU one. Name
    three things that changed, only one of which is ``two things run at once''.
  \item Why is \code{atomic_read} followed by \code{atomic_set} not an atomic read-modify-write?
  \item You hold a spinlock and call \code{kmalloc(GFP_KERNEL)}. What is wrong, and when will you
    find out?
  \item What exactly does a spinlock do to the CPU that is waiting for it? Why is that acceptable
    for short sections and catastrophic for long ones?
  \item Two functions take locks A and B in opposite orders, and the code has run in production for
    two years without deadlocking. Is it correct?
  \item Why does \code{CONFIG_PROVE_LOCKING} make a kernel slower, and why would you still enable it
    during development?
\end{enumerate}

\needspace{6\baselineskip}
\subsection*{Where to look}
\begin{itemize}
  \item \file{kernel/qa.config} turns on \code{CONFIG_PROVE_LOCKING} and
    \code{CONFIG_DEBUG_ATOMIC_SLEEP} for this course, with a comment on why. Those two options are
    the reason the exercises in this chapter produce output rather than silence.
\end{itemize}

\needspace{6\baselineskip}
\subsection*{Looking ahead}
Chapter~\ref{c8:ch} is about:
\begin{itemize}
  \item The other source of concurrency, and the one you cannot schedule.
  \item What an interrupt handler may not do, and what it must do before it returns.
  \item Four places to put work that does not belong in a handler.
  \item Why the default answer is now a thread.
\end{itemize}
